%!TEX encoding = UTF-8 Unicode
\documentclass[runningheads]{llncs}

%% ── packages ─────────────────────────────────────────────────────────────────
\usepackage{booktabs}
\usepackage{multirow}
\usepackage{makecell}
\usepackage{graphicx}
\usepackage{subcaption}
\usepackage{amsmath,amssymb,amsfonts}
\usepackage{algorithm}
\usepackage{algpseudocode}
\usepackage{listings}
\usepackage{xcolor}
\usepackage{microtype}
\usepackage{enumitem}
\setlist{nosep,leftmargin=*}
\usepackage{tikz}
\usepackage{pgfplots}
\pgfplotsset{compat=1.18}
\usepackage{hyperref}

%% ── reduce spacing ───────────────────────────────────────────────────────────
\setlength{\textfloatsep}{6pt plus 2pt minus 2pt}
\setlength{\floatsep}{4pt plus 2pt minus 2pt}
\setlength{\intextsep}{4pt plus 2pt minus 2pt}
\setlength{\abovecaptionskip}{3pt}
\setlength{\belowcaptionskip}{2pt}
\setlength{\parskip}{2pt plus 1pt minus 1pt}
\setlength{\parsep}{2pt plus 1pt minus 1pt}
\setlength{\itemsep}{1pt plus 1pt minus 1pt}

%% TikZ libraries for figures
\usetikzlibrary{shapes.geometric,arrows.meta,positioning,fit,backgrounds,calc}

%% ── listings style ───────────────────────────────────────────────────────────
\lstset{
  basicstyle=\ttfamily\scriptsize,
  breaklines=true,
  frame=single,
  numbers=left, numberstyle=\tiny,
  keywordstyle=\color{blue!70!black}\bfseries,
  commentstyle=\color{gray},
  stringstyle=\color{teal},
}

\begin{document}

\title{OASIS: Supplementary Materials}

\titlerunning{OASIS: Supplementary Materials}

\author{Author Names Redacted for Review}

\authorrunning{Author Names Redacted}

\institute{Affiliation Redacted\\
\email{redacted@review.org}}

\maketitle

This document collects supplementary materials for the paper
``OASIS: Feedback-Driven Statistics Correction for Cost-Based Query Optimizers.''
It provides detailed scope and error-budget positioning (\S\,A),
model architecture and instrumentation details (\S\,B),
baseline fairness discussion and QuickSel-H adaptation (\S\,C),
the portable statistics-format conversion interface with PostgreSQL MCV instantiation (\S\,D),
predicate-coverage sensitivity analysis (\S\,E),
observation aggregation study details (\S\,F),
optimizer-only EXPLAIN protocol details (\S\,G),
and detailed overhead breakdowns (\S\,H).

%% ═════════════════════════════════════════════════════════════════════════════
\section{Scope and Error Budget}
\label{sec:scope-app}
%% ═════════════════════════════════════════════════════════════════════════════

OASIS targets \emph{single-column histogram staleness}.
We explicitly exclude from scope:
\begin{itemize}
  \item \textbf{Multi-column correlation errors}: Join cardinality misestimation
        caused by correlated column predicates requires multi-dimensional
        statistics (e.g., mvdists, DNF-based sketches, or learned joint models).
        OASIS operates on one-dimensional marginals and is complementary to
        multi-column methods.
  \item \textbf{Physical design errors}: Missing indexes, suboptimal partitioning,
        or incorrect materialized views are outside the statistics layer.
  \item \textbf{Cost-model calibration}: Even with correct selectivities, the
        optimizer's cost model may mispredict operator runtimes.
        OASIS improves selectivity estimates but does not modify the cost model.
  \item \textbf{Data distribution shifts beyond drift}: Fundamental schema changes
        (column additions, type changes) or massive bulk loads that alter the
        distribution shape beyond what a compound drift operator captures are
        out of scope; these scenarios require a fresh \texttt{ANALYZE}.
\end{itemize}

\paragraph{Error budget decomposition.}
In a production optimizer, the total cardinality estimation error decomposes as:
\[
  \text{Total Error} = \underbrace{\text{Single-Column Error}}_{\text{OASIS target}}
  + \underbrace{\text{Correlation Error}}_{\text{multi-column methods}}
  + \underbrace{\text{Cost-Model Error}}_{\text{calibration}}
  + \underbrace{\text{Physical Design Error}}_{\text{indexing/partitioning}}.
\]
OASIS addresses the first term. Even perfect single-column statistics do not
eliminate the remaining terms, but accurate marginals are a necessary
prerequisite: learned multi-column estimators (e.g., copula-based models like
CoLSE~\cite{rathuwadu2025colse}) combine one-dimensional marginals with a
learned dependence structure, making single-column statistics quality a
foundational concern.

%% ═════════════════════════════════════════════════════════════════════════════
%% Include the existing appendix files
%% ═════════════════════════════════════════════════════════════════════════════

\section{Model and Instrumentation Details}
\label{sec:system-details-app}

This appendix collects implementation details trimmed from the main paper for
page budget reasons.

\subsection{Imputation-Layer Architecture}

The imputation layer (Stage~1) is a compact \emph{permutation-invariant set encoder} over
the stale-quantile and feedback tokens. It is built so the same model consumes any
feedback-window size and is invariant to observation order:

\begin{enumerate}
  \item \textbf{Tokenization}: the normalized stale quantile boundaries, lightweight
        metadata, and up to $K$ recent feedback observations (predicate type, boundary,
        estimated and actual selectivity) are embedded into a common token space, with a
        validity mask for unused slots.
  \item \textbf{Set encoding}: the feedback tokens are aggregated by a permutation-invariant
        encoder (self-attention over the unordered window followed by masked pooling), so the
        representation is invariant to observation order and handles a variable number of
        observations.
  \item \textbf{Residual head}: the pooled representation, fused with the encoded prior,
        predicts a correction delta $\boldsymbol{\delta}$ to the interior quantile
        boundaries; the output is clamped and monotonized to a valid CDF, so the model learns
        drift-induced deviations rather than absolute boundaries.
\end{enumerate}

\paragraph{Training objective.}
Unlike a reconstruction-trained prior, the model is trained on the composite
optimizer-facing objective of the main paper: the held-out future-predicate log-Q-error of
the \emph{projected} marginal, a FactorJoin bin-mass join-regret term, an in-window
feedback-consistency residual, and a light reconstruction regularizer toward the post-drift
marginal. Training uses a sparse-to-dense feedback curriculum and a domain-randomized mixture
of drift families (batch loads, range/date shifts, skew evolution, outlier bursts, multimodal
and seasonal drift) so the completion generalizes across feedback budgets and drift
mechanisms. The model is trained offline in PyTorch with Adam and reused across columns
without per-table retuning. The main-paper ablation shows that applying the projection inside
the training loop is inessential---post-hoc projection performs identically---so the
contribution is the training objective, not in-loop calibration.

\subsection{Per-Conjunct Operator Instrumentation}

Filter and scan-filter-project operators are augmented with per-conjunct
selectivity counters. During execution, whenever a batch of rows is processed
through a compound predicate $\phi_1 \wedge \phi_2 \wedge \cdots \wedge \phi_m$,
each conjunct $\phi_j$ is evaluated independently on the full input batch:
\[
  n_j^{\text{in}} \mathrel{+}= |\text{batch}|, \qquad
  n_j^{\text{out}} \mathrel{+}= |\{r \in \text{batch} : \phi_j(r)\}|.
\]
The final intersection of all conjunct results produces the same output as
the original compound filter while additionally recording per-conjunct
independent selectivities $s^*_j = n_j^{\text{out}} / n_j^{\text{in}}$.
Upon operator completion, these counters are written to the engine's existing
runtime metrics store. In the deployment model used by OASIS, the collector
simply piggybacks on predicate evaluation already being performed by the
engine, flushes a structured record to disk, and emits at most one
observation tuple per filter conjunct. As a result, each query writes at
most as many observation records as it has conjunctive filter clauses.

\begin{algorithm}[H]\small
\caption{Per-Conjunct Operator Instrumentation}
\label{alg:instrument}
\begin{algorithmic}[1]
\Require Compound filter $\Phi = \phi_1 \wedge \cdots \wedge \phi_m$, input batch $B$
\Ensure  Filtered output $B'$, counters $n_j^{\text{in}}, n_j^{\text{out}}$
\State $\text{mask} \leftarrow [\textbf{true}]^{|B|}$
\For{$j = 1$ \textbf{to} $m$}
  \State $\text{result}_j \leftarrow \phi_j(B)$
  \State $n_j^{\text{in}} \mathrel{+}= |B|$; \quad $n_j^{\text{out}} \mathrel{+}= |\text{result}_j|$
  \State $\text{mask} \leftarrow \text{mask} \wedge \text{result}_j$
\EndFor
\State \Return $B[\text{mask}]$
\Statex
\Statex \textit{// On operator completion, emit per-conjunct metrics:}
\For{$j = 1$ \textbf{to} $m$}
  \State \Call{EmitMetric}{conjunct\_$j$, $n_j^{\text{in}}$, $n_j^{\text{out}}$}
\EndFor
\end{algorithmic}
\end{algorithm}

\subsection{QuickSel-H Adaptation Details}

The QuickSel system~\cite{park2020quicksel} is a selectivity learner that fits a
mixture-model density to satisfy feedback constraints. Because QuickSel outputs
point selectivity estimates rather than histogram boundaries, we evaluate a
clearly labeled \textbf{QuickSel-H} adaptation that bridges the gap between
selectivity learning and histogram correction. The adaptation performs three steps:

\begin{enumerate}
  \item \textbf{Mixture-model fitting}: Given the same stale prior $H_0$ and
        observation window $\mathcal{O}$ of size $K{=}16$, QuickSel-H fits a
        Gaussian mixture model with $M{=}5$ components whose CDF is constrained
        to match observed selectivities at the predicate values in $\mathcal{O}$.
        The optimization uses the same relative-entropy objective as the original
        QuickSel, solved via iterated scaling.
  \item \textbf{Quantile extraction}: From the fitted mixture CDF, we extract
        the $B{-}1{=}9$ quantile values at the equi-depth levels used by all other
        methods. This ensures QuickSel-H outputs on the same histogram grid.
  \item \textbf{Output interface}: The extracted quantiles are passed through
        the same monotonic projection as OASIS and returned as the corrected histogram.
\end{enumerate}

\noindent\textbf{Why this adaptation is necessary.}
QuickSel's native output is a density estimate over the full domain, not a
fixed-resolution histogram consumable by the CBO. To compare under a matched
output interface, we must extract quantiles from this density. The extraction
step introduces a minor information bottleneck (compressing a rich density into
$B{-}1$ boundaries), but this is identical to the bottleneck all other methods
face and does not fundamentally change QuickSel's correction mechanism. We do
not claim QuickSel-H is a line-by-line reproduction of the original system; it
is a faithful adaptation that preserves QuickSel's core density-fitting approach
while conforming to the common output format required by our evaluation protocol.


\section{Baseline Fairness Discussion}
\label{sec:baseline-fairness}

All feedback-driven methods compared in the main paper consume the same
observation window of size $K{=}16$ and emit corrected quantiles on the same
$B{=}10$ equi-depth histogram grid. No method receives additional information
or computational budget beyond this shared interface.

\paragraph{STHoles.}
We implement the hierarchical refinement algorithm described
in~\cite{bruno2001stholes}. STHoles builds a tree of bucket refinements from
feedback. We initialize it from the stale prior and allow up to $K$ refinement
steps per correction instance, using the same observation tuples provided to
OASIS. The output quantiles are extracted from the refined tree on the standard
$B{-}1$ grid.

\paragraph{ISOMER.}
We implement the consistency-based histogram adjustment
from~\cite{markl2007consistent}. ISOMER adjusts bucket frequencies to satisfy
observed selectivity constraints while minimizing KL divergence from the prior.
We provide the same $K{=}16$ observation tuples and extract quantiles on the
standard grid.

\paragraph{QuickSel-H.}
See \S\,\ref{sec:system-details-app} for the detailed QuickSel-H adaptation.

\paragraph{Simple heuristics.}
LinInterp and FeedAvg consume the same observation window but use fixed
aggregation rules (blend factor 0.5 and exponential moving average,
respectively). They are included to demonstrate that learned correction is
necessary, not merely that any feedback-driven approach helps.

No method receives per-table or per-distribution tuning; all use default
hyperparameters. Results are averaged over three independent training seeds
(42, 123, 456) with 95\% confidence intervals.

\section{Portable Statistics-Format Conversion Interfaces}
\label{sec:mcv-adapter-app}

\subsection{Motivation}

OASIS corrects a canonical full-distribution view of a column. This is
straightforward when a DBMS exposes a self-contained histogram, but many
production systems do not. Instead, the optimizer's statistics object often
couples the histogram with auxiliary summaries whose probability mass must
remain consistent after any update.

\begin{itemize}
    \item \textbf{PostgreSQL}: Maintains a Most Common Values (MCV) list alongside a residual histogram that excludes MCV entries.
    \item \textbf{Oracle}: Uses top-frequency values together with height-balanced histograms.
    \item \textbf{SQL Server}: Employs a density vector together with histogram steps.
\end{itemize}

This coupling is the main portability obstacle for histogram-correction methods:
\emph{correcting the histogram alone} is often insufficient because the stored
statistics object is not a standalone histogram. OASIS addresses this with a
\emph{statistics-format conversion interface} that translates engine-specific
statistics layouts into a canonical full distribution, applies OASIS there, and
translates the corrected result back.

We instantiate the interface on PostgreSQL's MCV case because it is both
practically important and representative of tightly coupled statistics
architectures.

\subsection{Three-Phase Conversion Interface}

The interface follows three phases:

\begin{enumerate}
    \item \textbf{Reconstruction}: Convert engine-native statistics $\rightarrow$ canonical full distribution
    \item \textbf{Correction}: Apply OASIS on the canonical representation
    \item \textbf{Decomposition}: Convert corrected full distribution $\rightarrow$ engine-native statistics
\end{enumerate}

This design decouples OASIS's correction logic from database-specific storage
formats. PostgreSQL is the running example in this appendix, but the pattern is
more general.

\subsection{PostgreSQL Instantiation}

\subsubsection{Problem Formulation}

Given PostgreSQL statistics:
\begin{itemize}
    \item MCV list: $\mathcal{M} = \{(v_1, f_1), (v_2, f_2), \ldots, (v_k, f_k)\}$
    \item Residual histogram bounds: $\mathcal{H} = \{u_0, u_1, \ldots, u_m\}$ where adjacent values define $m$ equi-depth residual intervals
    \item Null fraction: $f_{\mathrm{null}}$
    \item Non-null invariant: $\sum_{i=1}^{k} f_i + p_{\mathrm{res}} = 1 - f_{\mathrm{null}}$
\end{itemize}

Construct a canonical full distribution $\mathcal{F}$ that exposes the complete
non-null column distribution to OASIS.

\subsubsection{Phase 1: Reconstruction}

\begin{algorithm}[H]
\caption{Reconstruct Canonical Full Distribution from PostgreSQL Statistics}
\label{alg:reconstruct}
\begin{algorithmic}[1]
\REQUIRE MCV list $\mathcal{M}$, residual bounds $\mathcal{H} = \{u_0, \ldots, u_m\}$, null fraction $f_{\mathrm{null}}$
\ENSURE Canonical full distribution $\mathcal{F} = (\mathcal{S}, \mathcal{R})$

\STATE Initialize $\mathcal{S} \leftarrow \emptyset$, $\mathcal{R} \leftarrow \emptyset$
\STATE $p_{\mathrm{res}} \leftarrow 1 - f_{\mathrm{null}} - \sum_{(v_i,f_i) \in \mathcal{M}} f_i$

\FOR{each $(v_i, f_i) \in \mathcal{M}$}
    \STATE Add singleton bucket $(v_i, f_i)$ to $\mathcal{S}$
\ENDFOR

\FOR{$j = 1$ \textbf{to} $m$}
    \STATE Add range bucket $(u_{j-1}, u_j, p_{\mathrm{res}} / m)$ to $\mathcal{R}$
\ENDFOR

\STATE Validate: $\sum_{s \in \mathcal{S}} s.\mathrm{freq} + \sum_{r \in \mathcal{R}} r.\mathrm{freq} = 1 - f_{\mathrm{null}}$
\RETURN $\mathcal{F} = (\mathcal{S}, \mathcal{R})$
\end{algorithmic}
\end{algorithm}

PostgreSQL stores histogram \emph{boundary values}, not explicit per-bucket
frequencies. The conversion interface therefore interprets the residual
histogram as $m$ equal-depth intervals carrying equal shares of the residual
non-MCV mass. Before tensorization, the interface evaluates the CDF of
$\mathcal{F}$ at OASIS's fixed quantile levels, yielding the canonical quantile
vector consumed by the model.

\subsubsection{Phase 2: OASIS Correction}

Once reconstructed, the canonical full distribution is passed unchanged to the
standard OASIS correction pipeline. The important systems property is that the
model itself remains ignorant of the source database format; all engine-specific
handling is localized to the conversion interface.

\subsubsection{Phase 3: Decomposition}

Given a corrected full distribution $\mathcal{F}' = (\mathcal{S}', \mathcal{R}')$,
reconstruct PostgreSQL-compatible statistics:
\begin{itemize}
    \item new MCV list: $\mathcal{M}'$
    \item new residual histogram bounds: $\mathcal{H}'$
    \item invariant: $\sum_i f_i' + p'_{\mathrm{res}} = 1 - f_{\mathrm{null}}$
\end{itemize}

\begin{algorithm}[H]
\caption{Decompose Canonical Full Distribution to PostgreSQL Statistics}
\label{alg:decompose}
\begin{algorithmic}[1]
\REQUIRE Full distribution $\mathcal{F}' = (\mathcal{S}', \mathcal{R}')$, MCV threshold $\tau$, MCV limit $K$, target histogram size $m{+}1$
\ENSURE MCV list $\mathcal{M}'$, residual histogram bounds $\mathcal{H}'$

\STATE \textbf{Step 1: Extract MCV candidates}
\STATE $\mathcal{C} \leftarrow \{(v, f) \mid (v, f) \in \mathcal{S}' \text{ and } f \geq \tau\}$
\STATE Sort $\mathcal{C}$ by frequency (descending)
\STATE $\mathcal{M}' \leftarrow$ top $K$ entries from $\mathcal{C}$

\STATE \textbf{Step 2: Build residual distribution}
\STATE $V_{\text{mcv}} \leftarrow \{v \mid (v, f) \in \mathcal{M}'\}$
\STATE Construct residual mixture from all range buckets in $\mathcal{R}'$ and all singletons in $\mathcal{S}'$ whose value is not in $V_{\text{mcv}}$

\STATE \textbf{Step 3: Resample equi-depth histogram bounds}
\FOR{$j = 0$ \textbf{to} $m$}
    \STATE Set $u_j$ to the residual quantile at level $j/m$
\ENDFOR
\STATE $\mathcal{H}' \leftarrow \{u_0, \ldots, u_m\}$

\STATE Validate: $\sum_{(v,f) \in \mathcal{M}'} f + p'_{\mathrm{res}} = 1 - f_{\mathrm{null}}$
\RETURN $\mathcal{M}'$, $\mathcal{H}'$
\end{algorithmic}
\end{algorithm}

\subsubsection{MCV Selection Policy}

The MCV threshold $\tau$ and limit $K$ control the final decomposition:
\begin{itemize}
    \item \textbf{High threshold}: keeps the MCV list compact but pushes more probability mass into the residual histogram
    \item \textbf{Low threshold}: retains more explicit singleton masses as MCV entries
    \item \textbf{PostgreSQL default-like regime}: $\tau \approx 0.01$, $K = 100$
\end{itemize}

This threshold is a decomposition policy choice, not part of OASIS's core
correction model. In our validation, experiments 1--3 remove threshold effects
by using threshold-free round trips, while experiment 4 isolates the effect of
$\tau_{\text{MCV}}$ after a simulated correction.

\subsection{Why This Is a Portable Interface, Not a PostgreSQL Patch}

The PostgreSQL instantiation is only one instance of the conversion pattern.
The same interface idea extends to other engines that rely on histograms to
very different degrees:

\begin{itemize}
    \item \textbf{Self-contained histograms} (MySQL, MariaDB): direct pass-through; conversion degenerates to a near-identity mapping.
    \item \textbf{MCV-coupled statistics} (PostgreSQL, Oracle): reconstruct from explicit frequent-value summaries plus residual histograms, then decompose back.
    \item \textbf{Hybrid summaries} (SQL Server): reconstruct from histogram steps plus density summaries, apply OASIS, then recompute auxiliary summaries.
\end{itemize}

The systems contribution is therefore broader than ``support PostgreSQL''.
The conversion interface turns histogram correction into a reusable component
for heterogeneous optimizer statistics architectures.

\subsection{Consistency Guarantees}

\subsubsection{Probability Mass Conservation}

\begin{theorem}[Probability Conservation]
If the source statistics satisfy the non-null invariant
$\sum_i f_i + p_{\mathrm{res}} = 1 - f_{\mathrm{null}}$, then:
\begin{enumerate}
    \item reconstruction produces a canonical full distribution with the same non-null mass;
    \item decomposition emits engine-native statistics whose MCV mass plus residual mass again equals $1 - f_{\mathrm{null}}$.
\end{enumerate}
\end{theorem}

\begin{proof}
Reconstruction assigns all explicit MCV mass to singleton buckets and all
remaining non-null mass to residual intervals by construction. Decomposition
selects an explicit MCV set and resamples the remainder as residual histogram
bounds, so the output MCV mass and residual mass partition the same corrected
non-null mass. Therefore the invariant is preserved.
\end{proof}

\subsubsection{Approximate Idempotence}

\begin{definition}[Approximate Idempotence]
Let $(\mathcal{M}', \mathcal{H}') = \mathrm{Decompose}(\mathrm{Reconstruct}(\mathcal{M}, \mathcal{H}), \tau, K)$. The conversion interface is $\epsilon$-idempotent if the round trip preserves total non-null mass within $\epsilon$ and preserves the residual quantile function within $\epsilon$ over a fixed evaluation grid.
\end{definition}

Perfect idempotence need not hold under thresholded decomposition because:
\begin{itemize}
    \item the MCV threshold may change which values remain explicit,
    \item low-frequency singleton mass may be folded back into the residual distribution,
    \item floating-point rounding accumulates.
\end{itemize}

Our threshold-free validation shows exact total-mass preservation and
machine-precision residual-quantile fidelity for the tested PostgreSQL-style
statistics.

\subsection{Implementation Interface}

We provide a Python reference interface that exposes the conversion pattern
explicitly:

\begin{lstlisting}[language=Python, caption=Statistics-Format Conversion Interface (PostgreSQL Instance)]
class PostgreSQLStatisticsAdapter:
    def to_full_histogram(
        self,
        pg_stats: PostgreSQLStatistics
    ) -> FullHistogram:
        """Phase 1: Reconstruction"""
        pass

    def from_full_histogram(
        self,
        full_hist: FullHistogram
    ) -> PostgreSQLStatistics:
        """Phase 3: Decomposition"""
        pass

    def round_trip_test(
        self,
        pg_stats: PostgreSQLStatistics
    ) -> Tuple[bool, str]:
        """Validate round-trip fidelity"""
        pass
\end{lstlisting}

The full reference interface is available in
\texttt{appendix/mcv\_adapter\_interface.py}. This appendix is the
supplementary counterpart of the compact portability discussion in the
main text and is intended to be kept numerically consistent with
\texttt{experiments/results/mcv\_validation\_results.json}.

\subsection{Performance Considerations}

\subsubsection{Time Complexity}

\begin{itemize}
    \item \textbf{Reconstruction}: $O(k + n)$ where $k = |\mathcal{M}|$ and $n = |\mathcal{H}|$
    \item \textbf{Decomposition}: dominated by MCV candidate sorting and residual quantile sampling
    \item \textbf{Total overhead}: negligible relative to correction and well within planning-time constraints in our PostgreSQL validation
\end{itemize}

\subsubsection{Space Complexity}

The canonical full distribution is temporary and small:
\[
\mathrm{Space}(\mathcal{F}) = \mathrm{Space}(\mathcal{M}) + \mathrm{Space}(\mathcal{H}) + O(1).
\]

In practice this overhead is acceptable because statistics objects are compact,
the canonical distribution exists only during correction, and no persistent
storage expansion is required.

\subsection{Validation Summary}
\label{sec:mcv-eval}

We validate the PostgreSQL instantiation of the conversion interface on
synthetic PostgreSQL-style statistics spanning uniform, normal, Zipfian, and
bimodal distributions.

\begin{table}[h]
\centering
\caption{PostgreSQL Statistics-Format Interface Validation Summary}
\label{tab:mcv-validation-summary}
\begin{tabular}{l l}
\toprule
\textbf{Metric} & \textbf{Result} \\
\midrule
Round-trip total mass error & $0$ \\
Residual quantile MAE (mean) & $2.4 \times 10^{-17}$ \\
Residual quantile MAE (max) & $9.2 \times 10^{-17}$ \\
Selectivity MAE & $1.1 \times 10^{-4}$ \\
Selectivity P99 & $3.8 \times 10^{-3}$ \\
Default-like overhead (100 MCV, 10--20 bins) & $0.066$--$0.073$ ms \\
Worst tested overhead (50 MCV, 50 bins) & $1.03$ ms \\
\bottomrule
\end{tabular}
\end{table}

These results support three claims. First, the conversion interface is exact in
mass preservation and essentially exact in residual-quantile reconstruction
under threshold-free round trips. Second, it introduces negligible estimation
bias and low latency on PostgreSQL-like settings. Third, threshold sensitivity
is graceful: increasing $\tau_{\text{MCV}}$ shrinks the MCV list and gradually
increases selectivity error, but does not compromise mass conservation.

\subsection{Takeaway}

The main value of this interface is not merely that it accommodates one
PostgreSQL-specific edge case. Rather, it makes histogram correction portable to
database systems whose optimizer statistics are stored as tightly coupled
multi-component objects. OASIS corrects the canonical full-distribution view;
the conversion interface is what makes that correction deployable across
heterogeneous DBMS statistics architectures.


%% ═════════════════════════════════════════════════════════════════════════════
\section{Predicate-Coverage Sensitivity Analysis}
\label{sec:coverage-sensitivity-app}
%% ═════════════════════════════════════════════════════════════════════════════

OASIS correction quality depends on the observation window covering the value
range relevant to future predicates. We conduct a controlled sensitivity
experiment at $q{=}10$ by restricting observations to a ``hot region'' of width
$w$ centered at the domain midpoint.

\begin{table}[h]
\centering
\caption{Predicate-Coverage Sensitivity at $q{=}10$}
\label{tab:coverage-sensitivity}
\begin{tabular}{c r r r}
\toprule
Coverage $w$ & Stale Prior Q-Error & OASIS Q-Error & Improvement \\
\midrule
1.00 (full) & 2.720 & \textbf{1.232} & 55\% \\
0.70        & 2.720 & \textbf{1.268} & 53\% \\
0.50        & 2.720 & \textbf{1.395} & 49\% \\
0.30        & 2.720 & \textbf{1.625} & 40\% \\
0.15        & 2.720 & \textbf{1.890} & 30\% \\
\bottomrule
\end{tabular}
\end{table}

At $w{=}0.7$ (70\% coverage), Q-Error degrades by only 3\% relative to full
coverage. At $w{=}0.3$, degradation reaches 40\% but OASIS still outperforms
the stale prior by 26\%. Even at $w{=}0.15$ (15\% coverage), OASIS provides
30\% improvement over the stale prior, demonstrating graceful degradation under
sparse feedback. The abstention policy (\S\,H in the main paper) provides a
natural guard for extreme sparsity scenarios.

%% ═════════════════════════════════════════════════════════════════════════════
\section{Observation Aggregation Study}
\label{sec:aggregation-app}
%% ═════════════════════════════════════════════════════════════════════════════

The main paper reports that replacing attention pooling with mean/max pooling
degrades Q-Error by 12\%/18\% at $q{=}10$. Here we provide the full breakdown.

\begin{table}[h]
\centering
\caption{Observation Aggregation Ablation at $q{=}10$}
\label{tab:aggregation}
\begin{tabular}{l r r}
\toprule
Pooling Strategy & Q-Error & Degradation vs.\ Attention \\
\midrule
Attention (OASIS) & \textbf{1.232} & --- \\
Mean pooling       & 1.380 & +12\% \\
Max pooling        & 1.453 & +18\% \\
\bottomrule
\end{tabular}
\end{table}

Mean pooling treats all observations equally, diluting the signal from
highly informative predicates. Max pooling amplifies the most extreme
observation but is sensitive to outliers. Attention learns non-uniform
weights that up-weight informative observations while down-weighting
redundant or noisy ones. This advantage is most pronounced when the
observation window contains heterogeneous predicates of varying
selectivity and informativeness.

%% ═════════════════════════════════════════════════════════════════════════════
\section{Optimizer-Only EXPLAIN Protocol Details}
\label{sec:optimizer-protocol-app}
%% ═════════════════════════════════════════════════════════════════════════════

The optimizer-only evaluation protocol uses PostgreSQL's \texttt{EXPLAIN}
(without \texttt{ANALYZE}) to validate the statistics$\rightarrow$plan causal
link. \texttt{EXPLAIN} invokes the full optimizer pipeline---parsing, rewriting,
planning, and cost estimation---but never executes the query, making it safe
for repeated invocation on production systems.

\paragraph{Protocol steps.}
\begin{enumerate}
  \item Create a synthetic table with 200K rows and known distribution.
  \item Run \texttt{ANALYZE} to capture fresh statistics.
  \item Apply controlled DML mutations to create drift.
  \item For each test query, run \texttt{EXPLAIN} under three statistics states:
        \textbf{(a)} stale (pre-mutation statistics),
        \textbf{(b)} fresh (post-mutation statistics captured by re-running
        \texttt{ANALYZE}),
        \textbf{(c)} OASIS-corrected (stale statistics with quantile boundaries
        replaced by OASIS predictions).
  \item Compare plan structures: scan types, join methods, join order.
\end{enumerate}

\paragraph{Results.}
On a TPC-DS-style scenario with extreme drift ($q \gg 10$):
\begin{itemize}
  \item 8 queries tested, 2/8 (25\%) showed scan-type transitions
        (Index Scan $\rightarrow$ Bitmap Heap Scan) when switching from stale
        to fresh statistics.
  \item OASIS-corrected statistics produced the same scan-type transitions
        as fresh statistics in both cases, confirming that the CBO's plan
        selection responds to the corrected quantile boundaries.
  \item Planning cost: $<5$\,ms per \texttt{EXPLAIN} invocation.
\end{itemize}

This protocol validates the causal chain
$\text{statistics correction} \rightarrow \text{selectivity change}
 \rightarrow \text{cost re-estimation} \rightarrow \text{plan change}$
without requiring query execution, making it suitable for production-safe
plan-impact validation.

%% ═════════════════════════════════════════════════════════════════════════════
\section{Detailed Overhead Breakdown}
\label{sec:overhead-app}
%% ═════════════════════════════════════════════════════════════════════════════

\begin{table}[h]
\centering
\caption{OASIS Runtime Overhead Breakdown (1000 trials)}
\label{tab:overhead}
\begin{tabular}{l r r r}
\toprule
Component & p50 (ms) & p95 (ms) & p99 (ms) \\
\midrule
Feature tensorization       & 0.38 & 0.42 & 0.45 \\
Model forward pass          & 0.66 & 0.72 & 0.78 \\
Total correction            & 1.04 & 1.18 & 1.25 \\
Statistics-format conversion & 0.07 & 0.07 & 0.07 \\
Combined correction + conversion & 1.11 & 1.25 & 1.32 \\
\midrule
Feedback collection (per tuple batch) & 0.12 & 0.18 & 0.23 \\
\bottomrule
\end{tabular}
\end{table}

For comparison, a column-level \texttt{ANALYZE} on a $10^7$-row table costs an
estimated 300--500\,ms of sequential I/O. OASIS's combined latency of
$\approx$1.13\,ms is ${\approx}300\times$ cheaper per invocation, well within
the 200\,ms planning budget (desideratum \textbf{D2}).

The statistics-format conversion interface adds negligible overhead
(0.066--0.073\,ms with PostgreSQL-like defaults of 100 MCV entries and 10--20
residual bins). Even in the worst tested configuration (50 MCV, 50 bins), the
conversion completes in 1.03\,ms.

\bibliographystyle{splncs04}
\bibliography{references}

\end{document}
